% PROPOSED replacement for sections/08_threats_limitations.tex. NOT APPLIED. States only limitations that hold regardless of the results.
\section{Threats to Validity and Limitations}
\label{sec:threats}
\textbf{Label validity.} The expected identity is derived from ORCID values, and OpenAlex profiles may carry ORCID values from overlapping upstream sources, so agreement with the benchmark is not independent verification. Independent verification covers every automatic assignment of the primary pair and an audit sample of its non-automatic cases. Results for the other cases, and all results for other conditions and policies, rest on benchmark-derived labels. The audited error rate describes only the sampled frame, is reported but not used for correction, and is not transferable. Uncertain and unverifiable cases are excluded from the primary rates and may concentrate among hard cases, so the primary rates may be optimistic; the bracket bounds this effect. An impure benchmark group is a finding about the benchmark and is reported separately.

\textbf{Population.} The frame holds benchmark-defined blocks with two to eight ORCID groups in the test split. Results do not extend to all author names, to the most ambiguous names excluded by the cap, to other splits, or to ambiguous queries for which no unique person exists; the benchmark and the smoke seeds contain none, and false acceptance on unresolved queries is not measured. The XL stratum has only nine test blocks, so it is described and not estimated.

\textbf{Review.} Reviewers can meet an ORCID in their own searches despite the blinding, and the independence of reviewers from the system and the benchmark labels rests on their declarations. Agreement statistics are unstable at small numbers. The Stage B evidence for a candidate is a capped list of works, which may not separate persons with common names.

\textbf{Sampling and inference.} Blocks are treated as independent after one case per block and distinct ORCIDs per block, but cases from different blocks may still share institutions, coauthor networks or candidate profiles, and the number of profiles that occur in more than one pool is reported as a diagnostic. Per-stratum binomial intervals are conservative approximations and not design-exact. In simulation, the pooled stratified bootstrap covered slightly below its nominal level (0.935), and the predictive sample-size rule fell about two percentage points short of its target when the coverage distribution differed from the prior and the pilot was small; these procedures are therefore accompanied by exact per-stratum intervals and by a conservative sensitivity calculation. The sample size is conditional on the coverage observed in a pilot of finite size.

\textbf{System.} The signal weights and thresholds are engineering defaults and are not validated or calibrated, and the score is a heuristic and not a probability. The coauthor signal cannot be observed on this benchmark, publication overlap is available only in the exploratory condition, and candidates scored on different signal sets are not comparable, which the policy treats by withholding automation. The designated citation of the primary context carries a weak hint of group membership. Affiliation cleaning is minimal and the normaliser is undecided. Candidates come from one provider whose ranking and content may drift over time; each capture is a hashed snapshot, and a later run may differ.

\textbf{Tooling.} The evaluation tooling was validated on synthetic data and a small smoke capture with references written by the authors, and the simulations assume independent strata, simple random sampling of blocks and Bernoulli outcomes. They show that the procedures behave as specified, not that the real data follow the simulated model.
